% GENERATED FILE. Edit canonical lessons, then run npm run generate:books.
% canonical-source-hash: fnv1a64:b23ed0d1b2757bc7
% canonical-lessons: HI-C360-apatti, HI-C360-ninda, HI-C360-drstikon, HI-C360-dhamkana, HI-C360-phuslana

\chapter{Objecting and Persuading}
\label{ch:hi-objecting-and-persuading}
\hlchaptermodality{360}

\begin{chapteropening}
\textbf{By the end of this chapter:} \emph{I can object, criticise, threaten or coax.}

Complete the last lesson of chapter 360: I can object, criticise, threaten or coax.
\end{chapteropening}

\section[\texorpdfstring{\hi{आपत्ति} (āpatti)}{आपत्ति (āpatti)}]{\hi{आपत्ति} (āpatti) — an objection}

\label{lesson:HI-C360-apatti}



\begin{quote}
Before the new one: say the Hindi for to get tangled, then the Hindi for to raise (a topic).
\end{quote}

\subsection*{You'll want to know: \hi{आपत्ति}}
\textbf{\hi{आपत्ति}} — \emph{āpatti} — \textquotedblleft{}an objection\textquotedblright{}.

\begin{quote}
\hi{मुझे} \hi{तुम्हारे} \hi{विचार} \hi{पर} \hi{कोई} \hi{आपत्ति} \hi{नहीं} \hi{है।}
\end{quote}

\begin{quote}
\emph{mujhe tumhāre vicār par koī āpatti nahīṃ hai.}
\end{quote}

\begin{quote}
I have no objection to your idea.
\end{quote}

Feminine.

\begin{cousinweb}
\textbf{\hi{आपत्ति}} (\emph{āpatti}) is a tatsama, borrowed from Sanskrit \emph{āpatti}, \textquotedblleft{}an occurrence; a mishap\textquotedblright{} (Monier-Williams). Monier-Williams does not give the sense of an objection, which is a later Hindi one.
\end{cousinweb}

\begin{grammarlens}[title={What you've built}]
One more word for giving reasons and opinions.
\end{grammarlens}

\subsection*{Your turn}
\begin{itemize}
\raggedright
  \item \emph{Say it:} \emph{āpatti}
  \item \emph{Say it:} \emph{āpatti}, once more
  \item \emph{Say it:} \emph{mujhe tumhāre vicār par koī āpatti nahīṃ hai.} — I have no objection to your idea.
  \item \emph{Recall:} say the Hindi for to get tangled, then the Hindi for to raise (a topic), then say \emph{āpatti} again
\end{itemize}

\begin{tcolorbox}[breakable,colback=teal!4,colframe=teal!35!black,title={Before you move on}]
What does \hi{आपत्ति} mean? (\textquotedblleft{}An objection\textquotedblright{}.) Say it once more.
\end{tcolorbox}
\section[\texorpdfstring{\hi{निंदा} (nindā)}{निंदा (nindā)}]{\hi{निंदा} (nindā) — condemnation, criticism}

\label{lesson:HI-C360-ninda}



\begin{quote}
Before the new one: say the Hindi for to raise (a topic), then the Hindi for an objection.
\end{quote}

\subsection*{You'll want to know: \hi{निंदा}}
\textbf{\hi{निंदा}} — \emph{nindā} — \textquotedblleft{}condemnation, criticism\textquotedblright{}.

\begin{quote}
\hi{सब} \hi{लोगों} \hi{ने} \hi{उसके} \hi{झूठ} \hi{की} \hi{निंदा} \hi{की।}
\end{quote}

\begin{quote}
\emph{sab logoṁ ne uske jhūṭh kī nindā kī.}
\end{quote}

\begin{quote}
Everyone condemned his lie.
\end{quote}

Feminine.

\begin{cousinweb}
\textbf{\hi{निंदा}} (\emph{nindā}) is a tatsama, borrowed from Sanskrit \emph{nindā}, \textquotedblleft{}blame, censure\textquotedblright{} (Monier-Williams).
\end{cousinweb}

\begin{grammarlens}[title={What you've built}]
One more word for giving reasons and opinions.
\end{grammarlens}

\subsection*{Your turn}
\begin{itemize}
\raggedright
  \item \emph{Say it:} \emph{nindā}
  \item \emph{Say it:} \emph{nindā}, once more
  \item \emph{Say it:} \emph{sab logoṁ ne uske jhūṭh kī nindā kī.} — Everyone condemned his lie.
  \item \emph{Recall:} say the Hindi for to raise (a topic), then the Hindi for an objection, then say \emph{nindā} again
\end{itemize}

\begin{tcolorbox}[breakable,colback=teal!4,colframe=teal!35!black,title={Before you move on}]
What does \hi{निंदा} mean? (\textquotedblleft{}Condemnation, criticism\textquotedblright{}.) Say it once more.
\end{tcolorbox}
\section[\texorpdfstring{\hi{दृष्टिकोण} (dṛṣṭikoṇ)}{दृष्टिकोण (dṛṣṭikoṇ)}]{\hi{दृष्टिकोण} (dṛṣṭikoṇ) — a point of view}

\label{lesson:HI-C360-drstikon}



\begin{quote}
Before the new one: say the Hindi for an objection, then the Hindi for condemnation.
\end{quote}

\subsection*{You'll want to know: \hi{दृष्टिकोण}}
\textbf{\hi{दृष्टिकोण}} — \emph{dṛṣṭikoṇ} — \textquotedblleft{}a point of view\textquotedblright{}.

\begin{quote}
\hi{हर} \hi{आदमी} \hi{का} \hi{अपना} \hi{दृष्टिकोण} \hi{होता} \hi{है।}
\end{quote}

\begin{quote}
\emph{har ādmī kā apnā dṛṣṭikoṇ hotā hai.}
\end{quote}

\begin{quote}
Everyone has their own point of view.
\end{quote}

Masculine.

\begin{cousinweb}
\textbf{\hi{दृष्टिकोण}} (\emph{dṛṣṭikoṇ}) joins two Sanskrit words, \emph{dṛṣṭi}, \textquotedblleft{}sight\textquotedblright{}, and \emph{koṇa}, \textquotedblleft{}an angle, a corner\textquotedblright{} (Monier-Williams): the angle one looks from. The compound is a modern one.
\end{cousinweb}

\begin{grammarlens}[title={What you've built}]
One more word for giving reasons and opinions.
\end{grammarlens}

\subsection*{Your turn}
\begin{itemize}
\raggedright
  \item \emph{Say it:} \emph{dṛṣṭikoṇ}
  \item \emph{Say it:} \emph{dṛṣṭikoṇ}, once more
  \item \emph{Say it:} \emph{har ādmī kā apnā dṛṣṭikoṇ hotā hai.} — Everyone has their own point of view.
  \item \emph{Recall:} say the Hindi for an objection, then the Hindi for condemnation, then say \emph{dṛṣṭikoṇ} again
\end{itemize}

\begin{tcolorbox}[breakable,colback=teal!4,colframe=teal!35!black,title={Before you move on}]
What does \hi{दृष्टिकोण} mean? (\textquotedblleft{}A point of view\textquotedblright{}.) Say it once more.
\end{tcolorbox}
\section[\texorpdfstring{\hi{धमकाना} (dhamkānā)}{धमकाना (dhamkānā)}]{\hi{धमकाना} (dhamkānā) — to threaten}

\label{lesson:HI-C360-dhamkana}



\begin{quote}
Before the new one: say the Hindi for condemnation, then the Hindi for a point of view.
\end{quote}

\subsection*{You'll want to know: \hi{धमकाना}}
\textbf{\hi{धमकाना}} — \emph{dhamkānā} — \textquotedblleft{}to threaten\textquotedblright{}.

\begin{quote}
\hi{दुकानदार} \hi{ने} \hi{चोर} \hi{को} \hi{धमकाया।}
\end{quote}

\begin{quote}
\emph{dukāndār ne cor ko dhamkāyā.}
\end{quote}

\begin{quote}
The shopkeeper threatened the thief.
\end{quote}

Past, with \hi{ने} (he or she …): \hi{उसने} \hi{मुझे} \hi{धमकाया}, \emph{usne mujhe dhamkāyā}.

\begin{cousinweb}
\textbf{\hi{धमकाना}} (\emph{dhamkānā}) is inherited from \emph{*dhammakka}, \textquotedblleft{}a threat\textquotedblright{}, a form CDIAL reconstructs from its descendants (CDIAL 6736).
\end{cousinweb}

\begin{grammarlens}[title={What you've built}]
One more word for giving reasons and opinions.
\end{grammarlens}

\subsection*{Your turn}
\begin{itemize}
\raggedright
  \item \emph{Say it:} \emph{dhamkānā}
  \item \emph{Say it:} \emph{dhamkānā}, once more
  \item \emph{Say it:} \emph{dukāndār ne cor ko dhamkāyā.} — The shopkeeper threatened the thief.
  \item \emph{Recall:} say the Hindi for condemnation, then the Hindi for a point of view, then say \emph{dhamkānā} again
\end{itemize}

\begin{tcolorbox}[breakable,colback=teal!4,colframe=teal!35!black,title={Before you move on}]
What does \hi{धमकाना} mean? (\textquotedblleft{}To threaten\textquotedblright{}.) Say it once more.
\end{tcolorbox}
\section[\texorpdfstring{\hi{फुसलाना} (phuslānā)}{फुसलाना (phuslānā)}]{\hi{फुसलाना} (phuslānā) — to coax, to wheedle}

\label{lesson:HI-C360-phuslana}



\begin{quote}
Before the new one: say the Hindi for a point of view, then the Hindi for to threaten.
\end{quote}

\subsection*{You'll want to know: \hi{फुसलाना}}
\textbf{\hi{फुसलाना}} — \emph{phuslānā} — \textquotedblleft{}to coax, to wheedle\textquotedblright{}.

\begin{quote}
\hi{बच्चा} \hi{दवा} \hi{नहीं} \hi{पी} \hi{रहा} \hi{था}, \hi{इसलिए} \hi{माँ} \hi{ने} \hi{उसे} \hi{फुसलाया।}
\end{quote}

\begin{quote}
\emph{baccā davā nahīṃ pī rahā thā, isliye māṁ ne use phuslāyā.}
\end{quote}

\begin{quote}
The child wouldn't take the medicine, so his mother coaxed him.
\end{quote}

Past, with \hi{ने} (he or she …): \hi{उसने} \hi{बच्चे} \hi{को} \hi{फुसलाया}, \emph{usne bacce ko phuslāyā}.

\begin{cousinweb}
\textbf{\hi{फुसलाना}} (\emph{phuslānā}) is a tadbhava, inherited from Sanskrit \emph{spṛśati}, \textquotedblleft{}touches\textquotedblright{} (CDIAL 13815). CDIAL lists the Hindi verb, \textquotedblleft{}to cajole\textquotedblright{}, under this entry; the step from touching to coaxing is not spelled out there.
\end{cousinweb}

\begin{grammarlens}[title={What you've built}]
One more word for giving reasons and opinions.
\end{grammarlens}

\subsection*{Your turn}
\begin{itemize}
\raggedright
  \item \emph{Say it:} \emph{phuslānā}
  \item \emph{Say it:} \emph{phuslānā}, once more
  \item \emph{Say it:} \emph{baccā davā nahīṃ pī rahā thā, isliye māṁ ne use phuslāyā.} — The child wouldn't take the medicine, so his mother coaxed him.
  \item \emph{Recall:} say the Hindi for a point of view, then the Hindi for to threaten, then say \emph{phuslānā} again
\end{itemize}

\begin{tcolorbox}[breakable,colback=teal!4,colframe=teal!35!black,title={Before you move on}]
What does \hi{फुसलाना} mean? (\textquotedblleft{}To coax, to wheedle\textquotedblright{}.) Say it once more.
\end{tcolorbox}
